\lecture{7.5}{Influence-Function-Based Inference}
\label{lec:7-5}

\subsection{Asymptotic Linearity}

An estimator $\psi_n$ is asymptotically linear for a parameter $\psi$ if it can be
written as
\[
\psi_n-\psi
=
\frac{1}{n}\sum_{i=1}^n D(O_i)+o_p(n^{-1/2}),
\]
where $D(O)$ has mean zero and finite variance. The function $D$ is an influence
function. This representation implies
\[
\sqrt n(\psi_n-\psi)\overset{d}{\longrightarrow}
N\{0,\Var(D(O))\}.
\]

For the counterfactual mean $\psi(1)=\E\{\Yone\}$, the efficient influence function
under the nonparametric point-treatment model is
\[
D_1(O)
=
\frac{\ind(A=1)}{g_1(W)}
\{Y-\Qbar(1,W)\}
+
\Qbar(1,W)-\psi(1).
\]
For $\psi(0)$,
\[
D_0(O)
=
\frac{\ind(A=0)}{g_0(W)}
\{Y-\Qbar(0,W)\}
+
\Qbar(0,W)-\psi(0).
\]
For the ATE, the influence function is
\[
D_{\ATE}(O)=D_1(O)-D_0(O).
\]

\subsection{Estimated Influence Function}

In practice, replace unknown quantities by estimates. For $\psi(1)$:
\[
D_{1,n}(O_i)
=
\frac{\ind(A_i=1)}{g_{n,1}(W_i)}
\{Y_i-\Qbar_n(1,W_i)\}
+
\Qbar_n(1,W_i)-\psi_{n,\AIPTW}(1).
\]
For $\psi(0)$:
\[
D_{0,n}(O_i)
=
\frac{\ind(A_i=0)}{g_{n,0}(W_i)}
\{Y_i-\Qbar_n(0,W_i)\}
+
\Qbar_n(0,W_i)-\psi_{n,\AIPTW}(0).
\]
For the ATE:
\[
D_{\ATE,n}(O_i)=D_{1,n}(O_i)-D_{0,n}(O_i).
\]

\subsection{Variance Estimation}

Estimate the asymptotic variance of $\psi_n(1)$ by the empirical variance of the
estimated influence function:
\[
\hat\tau_1^2
=
\frac{1}{n}\sum_{i=1}^n D_{1,n}(O_i)^2.
\]
The standard error is
\[
\widehat{\mathrm{SE}}\{\psi_n(1)\}
=
\sqrt{\hat\tau_1^2/n}.
\]
For the ATE, use
\[
\hat\tau_{\ATE}^2
=
\frac{1}{n}\sum_{i=1}^n
\left\{
D_{\ATE,n}(O_i)-\bar D_{\ATE,n}
\right\}^2,
\]
where $\bar D_{\ATE,n}=n^{-1}\sum_iD_{\ATE,n}(O_i)$, and
\[
\widehat{\mathrm{SE}}(\widehat{\ATE})
=
\sqrt{\hat\tau_{\ATE}^2/n}.
\]
For AIPTW, $\bar D_{\ATE,n}$ is usually very close to zero by construction, but
centering is harmless and often used in software.

\subsection{Wald Confidence Intervals}

An approximate 95\% Wald interval for $\psi(1)$ is
\[
\psi_{n,\AIPTW}(1)
\pm
1.96\sqrt{\hat\tau_1^2/n}.
\]
An approximate 95\% interval for the ATE is
\[
\widehat{\ATE}_{\AIPTW}
\pm
1.96\,\widehat{\mathrm{SE}}(\widehat{\ATE}).
\]

\begin{implementationbox}
Influence-function standard error for the ATE:
\begin{lstlisting}[language=R]
D1 <- (A == 1) / g_1W * (Y - Qbar_1W) + Qbar_1W - aiptw_ey1
D0 <- (A == 0) / g_0W * (Y - Qbar_0W) + Qbar_0W - aiptw_ey0
Date <- D1 - D0

se_ate <- sqrt(var(Date) / length(Y))
ci_ate <- aiptw_ate + c(-1.96, 1.96) * se_ate
\end{lstlisting}
\end{implementationbox}

\subsection{Hypothesis Tests}

For testing
\[
H_0:\ATE=0
\qquad\text{versus}\qquad
H_1:\ATE\neq 0,
\]
use
\[
Z=\frac{\widehat{\ATE}_{\AIPTW}}
{\widehat{\mathrm{SE}}(\widehat{\ATE})}.
\]
The two-sided Wald p-value is
\[
p=2\{1-\Phi(|Z|)\},
\]
where $\Phi$ is the standard normal distribution function.

\subsection{Coverage in Finite Samples}

Influence-function inference is asymptotic. In finite samples, coverage can be below
or above the nominal level if nuisance regressions are unstable, positivity is weak, or
the sample size is too small for the normal approximation. Coverage often improves as
$n$ grows, provided the nuisance estimators and regularity conditions are adequate.

The estimated influence-function values should be inspected. A few extremely large
values can indicate that confidence intervals are being driven by near-zero propensity
scores or extreme residuals.

\subsection{Simulation Interpretation}

Simulation studies often compare the empirical coverage probability of nominal 95\%
confidence intervals. If 1000 simulated datasets are generated and the 95\% interval
contains the true parameter in 946 of them, the estimated coverage probability is
94.6\%. Such studies reveal finite-sample behavior of the estimator and interval, not
whether a single real-data analysis satisfies the causal assumptions.

\begin{keyidea}
The influence function turns a complicated doubly robust estimator into an approximate
sample mean of mean-zero terms. This is the basis for standard errors, Wald intervals,
and p-values.
\end{keyidea}
